\scanpage{25}
(ii) $N'^*(X)$ est un \uline{idéal} de $A^*(X)$, donc $\operatorname{Cyc}^*(X)/N^*(X)\simeq A^*(X)/N'^*(X)$ est muni d'une structure d'\uline{anneau}, \unsure{noté} $A'^*(X)$

(iii) Pour tout hom $f:X\longrightarrow Y$ de \deleted{variétés projectives} \inserted{\unsure{qu. pr.}} \deleted{\unsure{fibrés}} lisses, l'hom $f^*:A^*(Y)\mathbin{-}A^*(X)$ applique $N'^*(Y)$ dans $N'^*(X)$, donc définit pr passage au quotient un hom d'anneaux $A'^*(Y)\mathbin{-}A'^*(X)$.

(iv) Pour tout $f$ comme dans (iii), \inserted{avec $f$ propre,} l'hom \inserted{de groupes} $f_*:A^*(X)\mathbin{-}A^*(Y)$ \deleted{est} applique $N'^*(X)$ dans $N'^*(Y)$, et induit donc un hom de groupes (également noté $f_*$)
\[
A'^*(X)\longrightarrow A'^*(Y)
\]

(v) L'équiv. d'Albanese $\Longrightarrow$ équivalence algébrique.

(Donc on peut parler d'éléments de $A'^*(X)$ \uline{alg. équiv. à $0$} !)

\uline{Dém.} (i) \deleted{Si $z$ est alb-équi. à $0$ sur $X$, il est évident qu'il en est de $=$ de $-z$.} \inserted{Évident $0$ est Alb-équi. à $0$.}

Prouvons que si $z,z'$ sont Alb-équiv. à $0$ sur $X$, il en est de $=$ de $z\unsure{-|+}z'$.

Soient $T,T'$ les \unsure{schémas} de paramètres comme dans 1.1. correspondant à $z,z'$. Soit $T_1=T\times T'$.

\mspage{24}
